% !TEX program = xelatex
% Based strictly on template.tex supplied by the user.

\documentclass[10pt,aspectratio=169]{beamer}

% ---------------- Packages ----------------
\usepackage[fontset=none]{ctex}
\usepackage{fontspec}
\usepackage{amsmath,amssymb,bm,mathtools}
\usepackage{booktabs,array,tabularx,multirow}
\usepackage{graphicx}
\usepackage{xcolor}
\usepackage{tikz}
\usepackage{adjustbox}
\usepackage{hyperref}

\usetikzlibrary{arrows.meta,positioning,calc}
\graphicspath{{figure/}}

% ---------------- Theme ----------------
\usetheme[numbering=fraction,progressbar=none,sectionpage=progressbar]{metropolis}
\metroset{
  titleformat=regular,
  titleformat title=regular,
  titleformat subtitle=regular,
  titleformat section=regular,
  titleformat frame=regular
}

% ---------------- Fonts ----------------
\IfFontExistsTF{Times New Roman}
  {\setmainfont{Times New Roman}\setsansfont{Times New Roman}}
  {\setmainfont{Noto Serif}\setsansfont{Noto Sans}}

\IfFontExistsTF{Kaiti SC}
  {\setCJKmainfont[AutoFakeBold=2.5]{Kaiti SC}
   \setCJKsansfont[AutoFakeBold=2.5]{Kaiti SC}}
  {\IfFontExistsTF{Noto Serif CJK SC}
    {\setCJKmainfont[AutoFakeBold=2.5]{Noto Serif CJK SC}
     \setCJKsansfont[AutoFakeBold=2.5]{Noto Serif CJK SC}}
    {\setCJKmainfont[
        Path=/usr/local/texlive/2024/texmf-dist/fonts/opentype/public/fandol/,
        BoldFont=FandolKai-Regular.otf
      ]{FandolKai-Regular.otf}
     \setCJKsansfont[
        Path=/usr/local/texlive/2024/texmf-dist/fonts/opentype/public/fandol/,
        BoldFont=FandolKai-Regular.otf
      ]{FandolKai-Regular.otf}}}

% ---------------- Colors ----------------
\definecolor{textred}{RGB}{125,34,51}
\definecolor{deepred}{RGB}{125,34,51}
\definecolor{darkred}{RGB}{92,12,20}
\definecolor{ink}{RGB}{20,20,20}
\definecolor{warmgray}{RGB}{92,88,84}
\definecolor{linegray}{RGB}{208,208,208}
\definecolor{footgray}{RGB}{228,228,228}
\definecolor{gold}{RGB}{180,130,40}

\setbeamercolor{palette primary}{fg=white,bg=deepred}
\setbeamercolor{title separator}{fg=textred}
\setbeamercolor{progress bar}{fg=deepred,bg=linegray}
\setbeamercolor{progress bar in section page}{fg=textred,bg=linegray}
\setbeamercolor{normal text}{fg=ink,bg=white}
\setbeamercolor{structure}{fg=deepred}
\setbeamercolor{alerted text}{fg=deepred}
\setbeamercolor{title}{fg=ink,bg=white}

% ---------------- Layout ----------------
\setbeamertemplate{navigation symbols}{}
\setbeamersize{text margin left=0.82cm,text margin right=0.82cm}
\setbeamerfont{normal text}{size=\fontsize{9.2}{13}\selectfont}
\setbeamerfont{frametitle}{size=\fontsize{15}{17}\selectfont,series=\mdseries}
\setbeamertemplate{itemize item}{\color{deepred}$\bullet$}
\setbeamertemplate{itemize subitem}{\color{textred}$\circ$}
\setlength{\tabcolsep}{3pt}
\linespread{1.02}

% ---------------- Frame title with progress line ----------------
\makeatletter
\setbeamertemplate{frametitle}{%
  \ifx\insertframetitle\@empty
  \else
    \nointerlineskip
    \vspace*{0.23cm}%
    \hspace*{-0.82cm}%
    \makebox[0pt][l]{\begin{tikzpicture}
      \pgfmathsetmacro{\progress}{\insertframenumber/\inserttotalframenumber}
      \node[
        anchor=west,
        inner sep=0pt,
        text=black,
        font=\fontsize{11}{11}\bfseries\selectfont
      ] at (0.75cm, 0.35cm) {\strut\insertframetitle};
      \draw[linegray, line width=0.6pt] (0, 0) -- (\paperwidth, 0);
      \draw[textred, line width=0.6pt] (0, 0) -- ({\progress*\paperwidth}, 0);
    \end{tikzpicture}}%
    \vspace{-0.02cm}
  \fi
}
\makeatother

% ---------------- Footer ----------------
\setbeamercolor{footleft}{fg=white,bg=darkred}
\setbeamercolor{footmid}{fg=deepred,bg=footgray}
\setbeamercolor{footright}{fg=deepred,bg=white}

\newcommand{\footerleft}{COPA-EEG}
\newcommand{\footermid}{Operator-Robust EEG Foundation Model}
\newcommand{\footerright}{Research Proposal}

\setbeamertemplate{footline}{%
  \leavevmode%
  \hbox{%
    \begin{beamercolorbox}[wd=0.30\paperwidth,ht=0.30cm,dp=0.06cm,center]{footleft}
      \raisebox{0.03cm}{\fontsize{7}{8}\selectfont \footerleft}
    \end{beamercolorbox}%
    \begin{beamercolorbox}[wd=0.50\paperwidth,ht=0.30cm,dp=0.06cm,center]{footmid}
      \raisebox{0.03cm}{\fontsize{7}{8}\selectfont \footermid}
    \end{beamercolorbox}%
    \begin{beamercolorbox}[wd=0.20\paperwidth,ht=0.30cm,dp=0.06cm,center]{footright}
      \raisebox{0.03cm}{\fontsize{7}{8}\selectfont \footerright\quad \insertframenumber/\inserttotalframenumber}
    \end{beamercolorbox}%
  }%
}

% ---------------- Reusable commands ----------------
\newenvironment{tightitemize}{%
  \begin{itemize}
  \setlength{\itemsep}{0.12em}
  \setlength{\parskip}{0pt}
}{\end{itemize}}

\newcommand{\hl}[1]{\textcolor{deepred}{#1}}
\newcommand{\subhead}[1]{\vspace{0.25em}{\fontsize{9.2}{13}\selectfont\color{deepred}\bfseries #1}\par\vspace{0.08em}}
\newcommand{\smallgray}[1]{{\fontsize{9.2}{13}\selectfont\color{ink}#1}}
\newcommand{\fig}[2][]{\includegraphics[#1]{#2}}
\newcommand{\source}[1]{{\vspace{0.08em}\fontsize{6.5}{8}\selectfont\color{warmgray}Source: #1\par}}
\newcommand{\eqnote}[1]{\vspace{-0.10em}{\fontsize{9.2}{13}\selectfont\color{ink}#1\par\vspace{0.34em}}}
\newcommand{\formulanote}[1]{{\fontsize{9.2}{13}\selectfont\hspace{2em}#1\par\vspace{0.12em}}}

\newenvironment{slidebody}{%
  \fontsize{9.2}{13}\selectfont
  \setlength{\abovedisplayskip}{0.30em}%
  \setlength{\belowdisplayskip}{0.30em}%
  \setlength{\abovedisplayshortskip}{0.18em}%
  \setlength{\belowdisplayshortskip}{0.18em}%
}{ }

\newcommand{\op}{\mathrm{op}}
\newcommand{\eqv}{\mathrm{eq}}
\newcommand{\pred}{\mathrm{pred}}
\newcommand{\var}{\mathrm{var}}
\newcommand{\orth}{\mathrm{orth}}
\newcommand{\softmax}{\operatorname{Softmax}}
\newcommand{\sg}{\operatorname{sg}}
\newcommand{\hsic}{\operatorname{HSIC}}
\newcommand{\R}{\mathbb{R}}

% ---------------- Title information ----------------
\title{COPA-EEG}
\subtitle{\textbf{面向观测算子鲁棒性的 EEG 基础模型}}
\author{{\fontsize{9.2}{13}\selectfont 研究方案汇报}}
\date{}
\institute{}

% ---------------- Section transition page ----------------
\setbeamertemplate{section page}{
  \centering
  \begin{minipage}{22em}
    \centering
    \usebeamercolor[fg]{section title}
    \usebeamerfont{section title}
    \insertsectionhead\\[-1ex]
    \usebeamertemplate*{progress bar in section page}
    \par
  \end{minipage}
  \par
  \vspace{\baselineskip}
}

\AtBeginSection[]{%
  \begin{frame}[plain,c,noframenumbering]
    \sectionpage
  \end{frame}
}

\begin{document}

\begin{frame}[plain]
\vspace*{-0.50cm}
\titlepage
\end{frame}

\section{研究动机}

\begin{frame}[t]{格式兼容不等于观测鲁棒}
\vspace{0.10cm}
{\fontsize{9.2}{13}\selectfont
\begin{columns}[T,totalwidth=\textwidth]
\begin{column}{0.48\textwidth}
\subhead{现有进展：输入兼容}
\begin{tightitemize}
  \item 任意通道数、电极位置与记录长度
  \item 时空依赖、多尺度与拓扑无关表示
  \item 跨任务、跨条件和跨 setup 预训练
\end{tightitemize}
\end{column}
\begin{column}{0.48\textwidth}
\subhead{仍未显式控制：观测混杂}
\begin{tightitemize}
  \item \hl{\textbf{reference}} 改变通道线性关系
  \item \hl{\textbf{filter}} 改变频谱相关结构
  \item \hl{\textbf{device}} 改变噪声与频率响应
\end{tightitemize}
\end{column}
\end{columns}
\vspace{0.35cm}
\[
\boxed{\text{同一神经活动}\;\not\Rightarrow\;\text{相同 token affinity}}
\]
}
\end{frame}

\begin{frame}[t]{输入兼容之后，还需控制表征关系}
\vspace{0.10cm}
\begin{center}
\begin{minipage}{0.94\textwidth}
{\color{black}
\renewcommand{\arraystretch}{1.20}
\setlength{\tabcolsep}{6pt}
\fontsize{8.7}{11.5}\selectfont
\begin{tabularx}{\textwidth}{>{\bfseries}m{0.20\textwidth}X X}
\toprule[0.9pt]
方法 & 已推进的问题 & 仍可能保留的混杂 \\
\midrule
LaBraM & 大规模自监督预训练 & 数据集与设备身份 \\
CBraMod / CSBrain & 时空解耦与多尺度建模 & operator-induced affinity \\
REVE / LUNA & 任意 setup 与拓扑无关表示 & compatibility $\neq$ invariance \\
NeurIPT & 跨任务、跨条件、电极配置 & 神经变化与算子变化未显式分离 \\
\bottomrule[0.9pt]
\end{tabularx}}
\end{minipage}
\end{center}
\vspace{0.22cm}
\eqnote{\centering \hl{\textbf{COPA-EEG 的研究位置：}}从统一输入空间，推进到控制观测算子造成的表征偏差。}
\end{frame}

\begin{frame}[t]{观测模型：网络看到 \(X_o\)，而非 \(Z\)}
\vspace{0.10cm}
\begin{slidebody}
\subhead{EEG 观测过程}
\vspace{0.04cm}
\[
\boxed{
X_o
=
M_oR_oG_o\Bigl[H_o\ast(LZ)\Bigr]+\epsilon_o
}
\]
\begin{columns}[T,totalwidth=\textwidth]
\begin{column}{0.48\textwidth}
\begin{tightitemize}
  \item \(L\)：脑源到头皮电极的传导
  \item \(R_o\)：参考算子
  \item \(M_o\)：montage 与通道选择
\end{tightitemize}
\end{column}
\begin{column}{0.48\textwidth}
\begin{tightitemize}
  \item \(H_o\)：滤波器与设备频响
  \item \(G_o\)：增益与单位变换
  \item \(\epsilon_o\)：噪声与伪影
\end{tightitemize}
\end{column}
\end{columns}
\vspace{0.22cm}
\[
Z\xrightarrow{o_1}X_{o_1},\qquad
Z\xrightarrow{o_2}X_{o_2},\qquad
X_{o_1}\not\equiv X_{o_2}.
\]
\end{slidebody}
\end{frame}

\begin{frame}[t]{普通 attention 混合了神经关系与算子关系}
\vspace{0.10cm}
\begin{slidebody}
\subhead{标准 self-attention}
\vspace{0.04cm}
\[
Q=XW_Q,\quad K=XW_K,\quad V=XW_V,\qquad
A=\softmax\!\left(\frac{QK^\top}{\sqrt d}\right).
\]
\formulanote{\(QK^\top\) 决定 token 如何建立连接；reference、filter 与 device 可系统性改变连接。}
\vspace{0.20cm}
\[
S_{ij}(X_o)
=
\underbrace{S_{ij}^{\mathrm{neural}}}_{\text{目标}}
+
\underbrace{S_{ij}^{\mathrm{operator}}(o)}_{\text{混杂}}.
\]
\end{slidebody}
\end{frame}

\section{方法框架}

\begin{frame}[t]{COPA-EEG：视图与 attention 双重控制}
\vspace{0.10cm}
\begin{slidebody}
\begin{center}
\begin{tikzpicture}[
  node distance=0.55cm,
  box/.style={draw=linegray,rounded corners=2pt,thick,minimum width=2.05cm,
    minimum height=0.85cm,align=center},
  key/.style={draw=deepred,rounded corners=2pt,thick,minimum width=2.25cm,
    minimum height=0.85cm,align=center},
  arrow/.style={-{Latex[length=2.2mm]},thick,draw=warmgray}
]
  \node[box] (x) {同一记录 \(X\)};
  \node[key,right=of x] (cop) {COP\\算子配对预训练};
  \node[box,right=of cop] (tok) {共享 backbone\\EEG tokens};
  \node[key,right=of tok] (opa) {OPA\\算子部分注意力};
  \node[box,right=of opa] (z) {表征 \(z\)};
  \draw[arrow] (x) -- (cop);
  \draw[arrow] (cop) -- (tok);
  \draw[arrow] (tok) -- (opa);
  \draw[arrow] (opa) -- (z);
\end{tikzpicture}
\end{center}
\vspace{0.26cm}
\begin{columns}[T,totalwidth=\textwidth]
\begin{column}{0.48\textwidth}
\subhead{Counterfactual Operator Pretraining}
\begin{tightitemize}
  \item 同一记录构造 operator-paired views
  \item 区分等价算子与信息有损算子
\end{tightitemize}
\end{column}
\begin{column}{0.48\textwidth}
\subhead{Operator-Partial Attention}
\begin{tightitemize}
  \item 学习算子相关低秩子空间
  \item 软抑制 \(Q,K\) 中的相关方向
\end{tightitemize}
\end{column}
\end{columns}
\end{slidebody}
\end{frame}

\begin{frame}[t]{COP：构造 operator-paired views}
\vspace{0.10cm}
\begin{slidebody}
\[
X^{(1)}=\mathcal T_{o_1}(X),
\qquad
X^{(2)}=\mathcal T_{o_2}(X).
\]
\begin{columns}[T,totalwidth=\textwidth]
\begin{column}{0.48\textwidth}
\subhead{近似等价算子}
\begin{tightitemize}
  \item 合法重参考、通道置换
  \item 单位换算、无截断增益
\end{tightitemize}
\[
\mathcal L_{\eqv}
=
\bigl\lVert z(X)-z\!\left(\mathcal T_o(X)\right)\bigr\rVert_2^2
\]
\end{column}
\begin{column}{0.48\textwidth}
\subhead{信息有损算子}
\begin{tightitemize}
  \item 通道删除、降采样、限带
  \item 强滤波、局部信号损坏
\end{tightitemize}
\[
\text{只预测仍可识别的信息}
\]
\end{column}
\end{columns}
\vspace{0.18cm}
\eqnote{\centering 不能把所有增强都视为“完整语义不变”。}
\end{slidebody}
\end{frame}

\begin{frame}[t]{有损视图只预测仍可识别的信息}
\vspace{0.10cm}
\begin{slidebody}
\subhead{Operator-conditioned latent prediction}
\vspace{0.04cm}
\[
\boxed{
\mathcal L_{\pred}
=
w(X,o)
\left\|g_o(z_s)-\sg(z_t)\right\|_2^2
}
\]
\begin{columns}[T,totalwidth=\textwidth]
\begin{column}{0.48\textwidth}
\subhead{角色分工}
\begin{tightitemize}
  \item \(z_t\)：完整视图 teacher
  \item \(z_s\)：退化视图 student
  \item \(g_o\)：算子条件预测器
\end{tightitemize}
\end{column}
\begin{column}{0.48\textwidth}
\subhead{可信权重 \(w(X,o)\)}
\begin{tightitemize}
  \item 保留通道比例与有效带宽
  \item teacher uncertainty
\end{tightitemize}
\end{column}
\end{columns}
\vspace{0.20cm}
\eqnote{\centering 避免退化视图迫使完整视图丢失有效神经信息。}
\end{slidebody}
\end{frame}

\begin{frame}[t]{OPA：学习算子相关低秩子空间}
\vspace{0.10cm}
\begin{slidebody}
\subhead{Operator embedding and subspace}
\vspace{0.04cm}
\[
e_o=E_{\op}(o),\qquad
U_{\ell,h}(o)\in\R^{d_h\times r},\quad r\ll d_h,
\]
\[
P_{\ell,h}(o)=U_{\ell,h}(o)U_{\ell,h}(o)^\top.
\]
\begin{columns}[T,totalwidth=\textwidth]
\begin{column}{0.48\textwidth}
\subhead{输入元数据}
\begin{tightitemize}
  \item reference、channel mask、电极坐标
  \item sampling rate、filter cutoff
  \item unit、gain、device metadata
\end{tightitemize}
\end{column}
\begin{column}{0.48\textwidth}
\subhead{可检验建模假设}
\begin{tightitemize}
  \item 算子影响在特征空间中可压缩
  \item 不同层与 attention head 影响不同
\end{tightitemize}
\end{column}
\end{columns}
\end{slidebody}
\end{frame}

\begin{frame}[t]{OPA 用软投影抑制 \(Q,K\) 中的算子方向}
\vspace{0.10cm}
\begin{slidebody}
\[
\widetilde Q_{\ell,h}
=
Q_{\ell,h}
\bigl(I-\alpha_{\ell,h}(o)P_{\ell,h}(o)\bigr),
\]
\[
\widetilde K_{\ell,h}
=
K_{\ell,h}
\bigl(I-\alpha_{\ell,h}(o)P_{\ell,h}(o)\bigr),
\qquad 0\le \alpha_{\ell,h}(o)\le 1,
\]
\[
A_{\ell,h}^{\mathrm{OPA}}
=
\softmax\!\left(
\frac{\widetilde Q_{\ell,h}\widetilde K_{\ell,h}^\top}{\sqrt{d_h}}
\right).
\]
\vspace{0.22cm}
\eqnote{\centering
\(\alpha=0\)：普通 attention
\quad|\quad
\(0<\alpha<1\)：部分抑制
\quad|\quad
\(\alpha=1\)：完全投影}
\end{slidebody}
\end{frame}

\begin{frame}[t]{从局部关系修正到最终表征约束}
\vspace{0.10cm}
\begin{slidebody}
\begin{columns}[T,totalwidth=\textwidth]
\begin{column}{0.48\textwidth}
\subhead{局部修正：attention relation}
\[
QK^\top\longrightarrow\text{token affinity}
\]
reference 与 montage 最直接改变通道间关系。
\end{column}
\begin{column}{0.48\textwidth}
\subhead{全局约束：representation leakage}
\[
V,\quad\text{residual},\quad\text{MLP}
\]
其他路径仍可能重新引入 operator information。
\end{column}
\end{columns}
\vspace{0.30cm}
\[
\widehat o=C_{\op}(z),\qquad
\mathcal L_{\op}=\hsic(z,o)
\quad\text{或 adversarial classification}.
\]
\end{slidebody}
\end{frame}

\begin{frame}[t]{总体目标：跨视图、去泄漏、稳几何}
\vspace{0.10cm}
\begin{slidebody}
\[
\boxed{
\mathcal L
=
\mathcal L_{\pred}
+\lambda_{\eqv}\mathcal L_{\eqv}
+\lambda_{\op}\mathcal L_{\op}
+\lambda_{\var}\mathcal L_{\var}
+\lambda_{\orth}\mathcal L_{\orth}
}
\]
\vspace{0.25cm}
\begin{center}
\begin{minipage}{0.86\textwidth}
{\color{black}
\renewcommand{\arraystretch}{1.18}
\setlength{\tabcolsep}{7pt}
\fontsize{8.8}{11.8}\selectfont
\begin{tabularx}{\textwidth}{>{\centering\arraybackslash}m{0.22\textwidth}X}
\toprule[0.9pt]
\textbf{约束} & \textbf{作用} \\
\midrule
\(\mathcal L_{\pred},\mathcal L_{\eqv}\) & 可识别信息预测与等价视图稳定性 \\
\(\mathcal L_{\op}\) & 最终表征的 operator leakage 控制 \\
\(\mathcal L_{\var},\mathcal L_{\orth}\) & 防止坍塌并约束低秩基几何 \\
\bottomrule[0.9pt]
\end{tabularx}}
\end{minipage}
\end{center}
\end{slidebody}
\end{frame}

\section{验证方案}

\begin{frame}[t]{四个可证伪假设}
\vspace{0.10cm}
\begin{center}
\begin{minipage}{0.92\textwidth}
{\color{black}
\renewcommand{\arraystretch}{1.25}
\setlength{\tabcolsep}{7pt}
\fontsize{8.8}{11.8}\selectfont
\begin{tabularx}{\textwidth}{>{\centering\arraybackslash\color{deepred}\bfseries}m{0.10\textwidth}X}
\toprule[0.9pt]
H1 & 物理有效算子视图之间存在可识别的共享神经信息。 \\
H2 & reference、montage、filter、device 在 \(Q,K\) 空间形成可压缩方向。 \\
H3 & 抑制这些方向可降低 attention 对观测条件的依赖，并保留下游关系。 \\
H4 & \(Q,K\) 修正不足以保证最终表征稳定，必须约束其他信息路径。 \\
\bottomrule[0.9pt]
\end{tabularx}}
\end{minipage}
\end{center}
\vspace{0.22cm}
\eqnote{\centering 不声称恢复唯一真实脑源；只检验跨观测系统的可迁移表征。}
\end{frame}

\begin{frame}[t]{性能提升必须伴随算子泄漏下降}
\vspace{0.10cm}
{\fontsize{9.2}{13}\selectfont
\begin{columns}[T,totalwidth=\textwidth]
\begin{column}{0.31\textwidth}
\subhead{分布迁移}
\begin{tightitemize}
  \item 跨设备
  \item 跨 reference
  \item 跨 montage
  \item 跨数据集
\end{tightitemize}
\end{column}
\begin{column}{0.31\textwidth}
\subhead{表征诊断}
\begin{tightitemize}
  \item operator probe
  \item \(Q/K/V\) leakage
  \item attention map
  \item final representation
\end{tightitemize}
\end{column}
\begin{column}{0.31\textwidth}
\subhead{关键消融}
\begin{tightitemize}
  \item Q-only / K-only
  \item QK / QKV
  \item synthetic / real shift
  \item frozen / few-shot / full
\end{tightitemize}
\end{column}
\end{columns}
\vspace{0.38cm}
\[
\boxed{
\text{跨域任务性能}\uparrow
\quad\land\quad
\text{operator predictability}\downarrow
}
\]
}
\end{frame}

\section{结论}

\begin{frame}[t]{将观测系统从隐含噪声变成显式变量}
\vspace{0.10cm}
\subhead{核心结论}
{\fontsize{9.2}{13}\selectfont
\begin{tightitemize}
  \item \hl{\textbf{问题重定义：}}EEG 异构性不仅是输入格式问题，也是 token relation 的混杂问题。
  \item \hl{\textbf{预训练重定义：}}等价算子学习稳定性；有损算子只预测仍可识别的信息。
  \item \hl{\textbf{attention 修正：}}逐层、逐 head 的低秩软投影抑制 operator-induced affinity。
  \item \hl{\textbf{证据标准：}}任务迁移、表征泄漏与结构消融必须形成闭环。
\end{tightitemize}
\vspace{0.35cm}
\[
\boxed{
\text{目标不是适配更多格式，而是让神经关系更少依赖观测系统}
}
\]
}
\end{frame}

\begin{frame}[t]{参考工作}
\vspace{0.10cm}
{\fontsize{7.2}{9.2}\selectfont
\begin{enumerate}
  \setlength{\itemsep}{0.18em}
  \item Jiang, W.-B., Zhao, L., Lu, B.-L.
    \emph{Large Brain Model for Learning Generic Representations with Tremendous EEG Data in BCI}. ICLR, 2024.
  \item Wang, J., et al.
    \emph{CBraMod: A Criss-Cross Brain Foundation Model for EEG Decoding}. ICLR, 2025.
  \item El Ouahidi, Y., et al.
    \emph{REVE: A Foundation Model for EEG}. NeurIPS, 2025.
  \item Fang, Z., et al.
    \emph{NeurIPT: Foundation Model for Neural Interfaces}. NeurIPS, 2025.
  \item Döner, B., et al.
    \emph{LUNA: Efficient and Topology-Agnostic Foundation Model for EEG Signal Analysis}. NeurIPS, 2025.
  \item Zhou, Y., et al.
    \emph{CSBrain: A Cross-scale Spatiotemporal Brain Foundation Model for EEG Decoding}. NeurIPS, 2025.
\end{enumerate}
}
\source{COPA-EEG\_汇报概述.md}
\end{frame}

\end{document}
